% 日本語の文書向けのtemplate.
% make pdf TEX=hello_world (LuaLaTeX) でbuild.
\documentclass{ltjsarticle}
% フォントは原ノ味 (TeX Live 同梱)。macOS のヒラギノに依存しないので Docker でも同じ出力になる
\usepackage[haranoaji,jis2004]{luatexja-preset}
\usepackage{amsmath, amssymb}
\usepackage{hyperref}
\hypersetup{setpagesize=false,
            colorlinks=true,
            linkcolor=blue,
            citecolor=blue,
            }

\begin{document}
\title{Hello \LaTeX!!}
\author{Shun Makino}
\date{}

\section{Hello \LaTeX!}
Hello \LaTeX!

\section{日本語の文章}
このように, 日本語の文章を出力することができます. (Build手順を間違えなければ)

\section{吾輩は猫である}
文献~\cite{neko}は日本を代表する小説の一つである.
\begin{quote}
    吾輩は猫である。名前はまだない。
    どこで生れたかとんと見当けんとうがつかぬ。何でも薄暗いじめじめした所でニャーニャー泣いていた事だけは記憶している。
\end{quote}

\bibliographystyle{junsrt}
\bibliography{reference.bib}

\end{document}
